% ==========================================================
% 对应位置的大小关系与交换对称
% 难度：★★ 中档
% ==========================================================

\newcommand{\qProbPermSymmPosStem}{%
将六棵高矮不同的小树栽成前后两排，每排三棵。求后排每棵小树都比其对应的前排小树高的概率。
}

\newcommand{\qProbPermSymmPosAnswer}{%
$\frac{1}{8}$
}

\newcommand{\qProbPermSymmPosSolution}{%
六棵树进入六个不同位置的总排列数为 $6!$。

将三个前后对应位置分别看成三列。对于每一列，确定该列的两棵树后，它们在前后两个位置上共有两种排列：高树在后、矮树在前；高树在前、矮树在后。其中恰有一种满足“后排高于前排”。

三列内部可以分别交换，因此对于任意一种按列配对的结果，共有 $2^3$ 种前后方向，而只有一种使三列都满足要求。

所以所求概率为 $\frac{1}{2^3} = \frac{1}{8}$。

（也可以计算有利排列数：$\frac{6!}{2^3}=90$，概率为 $\frac{90}{6!} = \frac{1}{8}$。）
}

\newcommand{\qProbPermSymmPosTeachingNotes}{%
\textbf{【避坑与边界说明】}
\begin{itemize}
    \item 条件是“后排每棵树高于其对应的前排树”，不是“后排任意一棵都高于前排任意一棵”。
    \item 每一列的两棵树一旦选定，满足条件的前后位置只有一种，因此不应再乘以 $2$。
    \item 三列之间的位置有区别，不需要除以 $3!$。
\end{itemize}
}
